\chapter*{TD 6 - Optimalité}

\section*{Rappels}
La condition d’échange entre intégrale et dérivation est typiquement garantie par l’existence pour tout $\theta$ d’un voisinage ouvert $V$ tel que $T sup_{\theta^\prime V} \frac{\partial}{\partial \theta} f_{\theta^\prime}$ soit intégrable ; le théorème de dérivation des intégrales à paramètres assure alorsla que la condition est vérifiée.

Dans les TD, la borne de Cramer Rao et information de Fisher ne seront évoquées que si cette condition
est implicitement satisfaite.

\medskip

\setcounter{exercice}{0}

\begin{mdframed}
\begin{exercice}
Soit $X_1, \dots, X_n$ un $n$-échantillon de densité 

$$
f(x) = \frac{1}{\theta} \exp\left(-\frac{x}{\theta}\right) \mathbb{1}_{[0,+\infty[}(x)
$$

\begin{enumerate}
    \item Montrer qe l'estimateur du maximum de vraisemblance de $\theta$ est $\hat{\theta} = \bar{X}$
    \item Déterminer l'information de Fisher du modèle.
    \item L'estimateur $\hat{\theta}$ est-il efficace pour estimer $\theta$ ?
\end{enumerate}
\end{exercice}
\end{mdframed}

\medskip

\begin{mdframed}
\begin{exercice}
Soit $(X_1, \dots, X_n)$ un $n$-échantillon de densité 

$$
f(x) = 2\theta x e^{-\theta x^2} \mathbb{1}_{x>0},
$$

où $\theta > 0$. On pourra admettre que si $Z \sim \Gamma(\alpha, \beta)$, alors :

\begin{itemize}
\item $\mathbb{E}[Z] = \frac{\alpha}{\beta}$, 
\item $\mathbb{E}[Z^2] = \frac{\alpha(\alpha+1)}{\beta^2}$, 
\item $\mathbb{E}\left[\frac{1}{Z}\right] = \frac{\beta}{\alpha-1}$, 
\item $\mathbb{E}\left[\frac{1}{Z^2}\right] = \frac{\beta^2}{(\alpha-1)(\alpha-2)}$
\end{itemize}

(propriétés à savoir retrouver).
\begin{enumerate}
    \item Calculer l'estimateur du maximum de vraisemblance $\hat{\theta}$ de $\theta$.
    \item Montrer que $Y_i = \theta X_i^2$ suit la loi exponentielle de paramètre $1$. En déduire la loi de $Z = \sum_{i=1}^n \theta X_i^2$.
\item Retrouver les propriétés de $Z$.
\item Déterminer l'information de Fisher du modèle.
    \item L'EMV est-il biaisé ? En déduire un estimateur $\hat{\theta}_1$ non biaisé.
    \item Calculer la variance de $\hat{\theta}_1$. L'estimateur $\hat{\theta}_1$ est-il efficace pour estimer $\theta$ ?
\end{enumerate}
\end{exercice}
\end{mdframed}

\medskip

\begin{mdframed}
\begin{exercice}
Soit $(X_1, \dots, X_n)$ $n$ variables i.i.d. suivant une loi de Poisson de paramètre $\theta$. On rappelle les propriétés suivantes de cette loi :

\begin{itemize}
    \item Sa fonction génératrice des moments est $M(t) = \exp(\theta(e^t - 1))$ pour $t \in \mathbb{R}$.
    \item $\mathbb{E}[X] = \theta$ ; $\mathbb{E}[X^2] = \theta^2 + \theta$ ; $\mathbb{E}[X^3] = \theta^3 + 3\theta^2 + \theta$ ; $\mathbb{E}[X^4] = \theta^4 + 6\theta^3 + 7\theta^2 + \theta$.
    \item $\sum_{i=1}^n X_i$ suit la loi de Poisson de paramètre $n\theta$.
\end{itemize}

On rappelle également :

\begin{itemize}
    \item L'estimateur du maximum de vraisemblance de $\theta$ est la moyenne empirique $\bar{X}$.
    \item L'information de Fisher du modèle vaut $I_n(\theta) = \frac{n}{\theta}$.
    \item L'estimateur $\bar{X}$ atteint la borne de Cramér--Rao : $\text{Var}(\bar{X}) = \frac{\theta}{n} = \frac{1}{I_n(\theta)}$ et est donc efficace pour estimer $\theta$.
\end{itemize}

\begin{enumerate}
    \item Dans cette question, on cherche à estimer $g(\theta) = e^{-\theta}$. On pose $Y_i = \mathbb{1}_{\{X_i=0\}}$ et $\bar{Y} = \frac{1}{n}\sum_{i=1}^n Y_i$.
    \begin{enumerate}
        \item Vérifier que $\bar{Y}$ est un estimateur sans biais de $g(\theta)$.
        \item Calculer la borne de Cramér--Rao pour $\bar{Y}$.
        \item Cet estimateur est-il efficace pour estimer $g(\theta)$ ?
    \end{enumerate}
    \item Dans cette question, on cherche à estimer $h(\theta) = \theta^2$. On pose $T_n = \bar{X}^2 - \frac{1}{n}\bar{X}$.
    \begin{enumerate}
        \item Vérifiez que $T_n$ est un estimateur sans biais de $h(\theta)$.
        \item Calculer la borne de Cramér--Rao pour $T_n$.
        \item Cet estimateur est-il efficace pour estimer $h(\theta)$ ?
    \end{enumerate}
\end{enumerate}
\end{exercice}
\end{mdframed}






























\medskip

\begin{mdframed}
\begin{exercice}
Soit $\theta > 0$ et $(X_1, \dots, X_n)$ un $n$-échantillon de densité :

$$
f_\theta(x) = \begin{cases} 
\frac{1}{\theta} x^{\frac{1}{\theta} - 1} & \text{si } 0 < x < 1, \\ 
0 & \text{sinon.} 
\end{cases} 
$$

\begin{enumerate}
    \item Calculer l'estimateur du maximum de vraisemblance $\hat{\theta}$ de $\theta$.
    \item Cet estimateur est-il biaisé ? \\
    \textit{Indication :} déterminer la loi de probabilité de $Y_i = -\frac{1}{\theta} \log X_i$ puis de $Z = \sum_{i=1}^n Y_i$.
    \item Calculer la variance de $\hat{\theta}$.
    \item Montrer que $\ell'(\theta) = \frac{n}{\theta^2}(\hat{\theta} - \theta)$, où $\ell(\theta)$ désigne la log-vraisemblance des observations.
    \item En déduire que l'estimateur $\hat{\theta}$ atteint la borne de Cramér-Rao.
\end{enumerate}
\end{exercice}
\end{mdframed}

\medskip

\begin{mdframed}
\begin{exercice}
On considère le modèle des lois $P_\theta$ de densité $f_\theta$, pour $\theta > 0$, données par :

$$
f_\theta(x) = \frac{\theta \exp\left(-\theta(\sqrt{x} - 1)\right)}{2\sqrt{x}} \mathbb{1}_{x \ge 1} 
$$

\begin{enumerate}
    \item Si $X$ est de loi $P_\theta$, quelle est la loi de $\sqrt{X}$ ?
    \item Vérifier que la variable aléatoire $\frac{\partial \ln f_\theta(X)}{\partial \theta}$ est centrée sous $P_\theta$.
    \item Calculer l'information de Fisher du modèle.
    \item On considère un $n$-échantillon $X_1, \dots, X_n$ de loi $P_\theta$. Montrer qu'il existe un unique estimateur du maximum de vraisemblance $\hat{\theta}_n$ et le calculer.
    \item Soit $g$ la fonction donnée par $g(u) = \frac{1}{u}$ pour tout $u > 0$. Écrire l'inégalité de Cramér-Rao pour l'estimation de $g(\theta)$.
    \item Montrer qu'il existe un estimateur sans biais efficace de $g(\theta)$.
\end{enumerate}
\end{exercice}
\end{mdframed}